% ============================================================
%  Наративний огляд літератури
%  Тема: Методика організації групової взаємодії учнів
%        у процесі пізнавальної діяльності
%        на заняттях «Я досліджую світ»
% ============================================================
%  Компілятор: XeLaTeX (рекомендовано) або LuaLaTeX
%  Бібліографія: biber  (не bibtex!)
%  Команда збірки:
%    xelatex review_template
%    biber   review_template
%    xelatex review_template
%    xelatex review_template
% ============================================================

\documentclass[a4paper,12pt]{article}

% ── Кодування та мова ──────────────────────────────────────
\usepackage{fontspec}           % XeLaTeX: довільні шрифти
\usepackage{polyglossia}
\setdefaultlanguage{ukrainian}
\setotherlanguage{english}
% Шрифт з підтримкою кирилиці (Times New Roman або Liberation Serif)
% На власному ПК / Overleaf: Times New Roman
\setmainfont{Times New Roman}
% Якщо Times New Roman відсутній (Linux-сервер), замініть на:
% \setmainfont{FreeSerif}  або  \setmainfont{DejaVu Serif}

% ── Поля сторінки ──────────────────────────────────────────
\usepackage[
  top    = 2.5cm,
  bottom = 2.5cm,
  left   = 3.0cm,
  right  = 1.5cm
]{geometry}

% ── Міжрядковий інтервал ───────────────────────────────────
\usepackage{setspace}
\onehalfspacing                 % полуторний інтервал

% ── Відступ першого рядка абзацу ───────────────────────────
\usepackage{indentfirst}
\setlength{\parindent}{1.25cm}

% ── Цитування (APA-подібний стиль) ─────────────────────────
\usepackage[
  backend  = biber,
  style    = apa,
  sorting  = nyt,               % за прізвищем → роком → назвою
  maxnames = 3,
  minnames = 1,
  uniquelist = false
]{biblatex}
\usepackage{csquotes}           % коректні лапки для biblatex
\addbibresource{references.bib}

% ── Гіперпосилання ─────────────────────────────────────────
\usepackage[
  colorlinks = true,
  linkcolor  = black,
  citecolor  = black,
  urlcolor   = blue,
  unicode    = true
]{hyperref}

% ── Додаткові пакети ───────────────────────────────────────
\usepackage{microtype}          % поліпшення мікротипографіки
\usepackage{booktabs}           % таблиці (за потреби)
\usepackage{enumitem}           % налаштування списків
\usepackage{xcolor}             % кольори (для TODO-нотаток)

% ── Допоміжна команда для TODO-приміток ────────────────────
\newcommand{\todo}[1]{\textcolor{red}{\textbf{[TODO: #1]}}}

% ── Метадані документа ─────────────────────────────────────
\title{%
  \large НАЗВА УНІВЕРСИТЕТУ\\[0.5em]
  \large Кафедра початкової освіти\\[2em]
  \LARGE\bfseries
  Методика організації групової взаємодії учнів\\
  у процесі пізнавальної діяльності\\
  на заняттях «Я досліджую світ»\\[1em]
  \large Наративний огляд літератури\\[2em]
}
\author{%
  Прізвище Ім'я По-батькові\\
  \small Спеціальність 013 Початкова освіта\\
  \small Науковий керівник: \underline{\hspace{6cm}}, к.пед.н., доцент
}
\date{\the\year}

% ============================================================

\begin{document}

\maketitle
\thispagestyle{empty}
\newpage

\tableofcontents
\newpage

% ==========================================================
% МЕТОДОЛОГІЧНА ПРИМІТКА
% ==========================================================

\section*{Методологічна примітка}
\addcontentsline{toc}{section}{Методологічна примітка}

Відповідно до статті~8, частини~4 Закону України «Про освіту» та принципів академічної доброчесності, задекларованих в університеті, автор зазначає, що у процесі підготовки цього огляду штучний інтелект (систему Claude компанії Anthropic) було використано виключно для технічної підтримки: первинної фільтрації та сортування бібліографічних записів із бази даних Scopus за формалізованими критеріями релевантності, тематичного групування відібраних джерел та форматування BibTeX-файлу і LaTeX-документа. Усі аналітичні судження, інтерпретація змісту джерел, перефразування та синтез ідей є результатом самостійної інтелектуальної роботи автора. Використання інструментів ШІ для технічної підтримки дослідницького процесу не є порушенням академічної доброчесності за умови його прозорого розкриття, що цим і здійснюється.

\newpage

% ==========================================================
% ВСТУП
% ==========================================================

\section{Вступ}

Групова взаємодія учнів у пізнавальній діяльності є одним із ключових механізмів реалізації компетентнісного підходу в початковій школі. В умовах Нової української школи інтегрований курс «Я досліджую світ» створює унікальний контекст для організації спільної дослідницької діяльності молодших школярів: поєднання природничого, соціального та громадянознавчого змістів у єдиному навчальному просторі органічно передбачає роботу у групах, що досліджують, обговорюють і пояснюють явища навколишнього світу. Незважаючи на значний масив міжнародних досліджень з кооперативного та inquiry-based навчання, питання методики організації саме групової взаємодії в контексті інтегрованих дослідницьких курсів початкової школи залишається недостатньо систематизованим у вітчизняній педагогічній науці.

Мета цього огляду — систематизувати міжнародні наукові дані щодо психолого-педагогічних засад, типів, вікових особливостей та методичних підходів до організації групової взаємодії учнів початкової школи у процесі пізнавальної діяльності дослідницького характеру. Для реалізації мети проведено наративний огляд літератури: пошук джерел здійснювався у базі даних Scopus за тринадцятьма тематичними запитами (блоки: кооперативне навчання, inquiry-based learning, інтегровані курси, роль учителя, оцінювання), хронологічний діапазон — 2015--2026 роки. За результатами первинної фільтрації (280 записів) і відбору за критеріями релевантності до огляду включено 63~джерела, з яких безпосередньо процитовано 45.

\newpage

% ==========================================================
% РОЗДІЛ 2 — ТЕОРЕТИЧНІ ЗАСАДИ
% ==========================================================

\section{Теоретичні засади групової взаємодії учнів}

\subsection{Психолого-педагогічні основи кооперативного навчання}

Кооперативне навчання є однією з найбільш доказово обґрунтованих стратегій організації навчального процесу в початковій школі. Систематичний мета-аналіз, що охоплює тридцять емпіричних досліджень, підтвердив статистично значущий позитивний вплив кооперативного навчання на математичну успішність учнів початкової школи, наголошуючи водночас на тому, що результативність стратегії залежить від якості її методичного впровадження \parencite{Talkhan2025}. Теоретичним підґрунтям цього підходу слугує ідея про те, що спільне розв'язання задач активізує когнітивні й метакогнітивні процеси, недосяжні в умовах індивідуальної роботи: учні вербалізують міркування, отримують зворотний зв'язок від однолітків і коригують власні способи дій \parencite{Gu2015143}.

Окрім академічної успішності, кооперативне навчання справляє доведений вплив на соціально-емоційний розвиток учнів. Квазіекспериментальне дослідження, проведене в початкових школах Італії за участі тридцяти шести класів, показало, що структуроване кооперативне навчання з технологічною підтримкою суттєво розвиває соціальні й емоційні навички молодших школярів порівняно з контрольною групою \parencite{Zagni2025}. Довгострокове лонгітюдне спостереження учнів п'ятого і шостого класів підтвердило, що кооперативна методологія асоціюється зі стабільнішим розвитком творчого мислення та лексичного збагачення протягом усього пізнього дитинства \parencite{Segundo-Marcos2025}.

Ефективне впровадження кооперативного навчання потребує цілеспрямованої підготовки вчителя і врахування специфіки класного середовища. Якісне дослідження на основі глибинних інтерв'ю з тридцятьма двома вчителями початкової школи в Ірані виявило три ключові групи перешкод: організаційні, педагогічні та пов'язані з учнівською поведінкою, і запропонувало конкретні стратегії їх подолання \parencite{Keramati2025120}. Пілотна програма в гетерогенних класах Женеви продемонструвала, що правильно організована кооперативна діяльність захищає соціальний клімат класу і знижує рівень ізоляції окремих учнів \parencite{Buchs202568}. Водночас аналіз динаміки групової колаборації у четвертому класі виявив, що позаурочні взаємодії між учнями виконують регулятивну функцію в структурі групової роботи, перерозподіляючи владні відносини всередині команди \parencite{Langer-Osuna2020514}. Ці дані вказують на те, що ефективність кооперативного навчання не є автоматичним наслідком групування учнів, а залежить від усвідомленого педагогічного управління груповою взаємодією.

% ============================================================
%  Розділ 2.2 — Типи групової роботи
%  ~300 слів
% ============================================================

\subsection{Типи групової роботи}

У сучасній педагогіці розрізняють кілька форм організації спільної навчальної діяльності, що різняться за ступенем взаємозалежності учасників, характером завдання і рівнем автономії групи. Кооперативне навчання передбачає чіткий розподіл ролей та спільну відповідальність за результат, тоді як колаборативне — більш гнучку, рівноправну взаємодію, орієнтовану на спільне конструювання смислів. Дослідження поведінки першокласників (7 років) у малих групах із трьох осіб за методологією Lesson Study засвідчило, що навіть наймолодші учні здатні до продуктивної кооперації за умови відповідного педагогічного супроводу, спростовуючи поширене переконання про «незрілість» дітей цього віку для групової роботи \parencite{Jakavonytė-Staškuvienė2025}. Дослідження спільного сторітелінгу в учнів 1--3 класів (97 осіб) показало, що тангібельні цифрові форми колаборації стимулюють більш інклюзивні та просоціальні взаємодії порівняно з нецифровими форматами, а колаборативні оповіді відрізняються вищою змістовою і формальною якістю \parencite{Filosofi2025801}.

Проєктна робота як окремий тип групової навчальної діяльності поєднує ознаки кооперації, дослідництва і продуктивної діяльності. Аналіз педагогіко-технологічних характеристик платформ електронного навчання для початкової школи виявив, що найефективніші з них підтримують усі фази проєктної роботи: генерацію ідей, розподіл завдань, спільне виконання та зворотний зв'язок \parencite{Baziukė2025}. Дослідницька групова діяльність у природничій освіті початкової школи набуває форми запитання-орієнтованого навчання, де учні колективно генерують гіпотези та обговорюють результати спостережень \parencite{Sakamoto2016105}.

Вибір типу групової роботи суттєво залежить від фізичного й технологічного середовища навчання. Дослідження у мультисенсорному класі з мульти-тач столами показало, що просторове розташування груп впливає на характер взаємодії учнів і розподіл навчальних функцій між членами групи \parencite{Mercier2016504}. Порівняльне дослідження ефектів кооперативного навчання в синхронному онлайн-форматі та в класі підтвердило, що обидва формати є результативними, однак відрізняються за характером міжособистісної взаємодії та груповою динамікою \parencite{Ralević2025}. Аналіз сприйняття кооперативного навчання учнями сільських початкових шкіл Китаю зафіксував, що більшість учнів оцінює групову роботу позитивно, пов'язуючи її з покращенням академічних результатів і розвитком комунікативних умінь \parencite{Xiao2025117}. Наведені дані свідчать про необхідність диференційованого підходу до вибору типу групової роботи з урахуванням навчального контексту, вікових особливостей учнів та дидактичної мети заняття.

\subsection{Вікові особливості групової взаємодії молодших школярів}

Молодший шкільний вік охоплює період від 6--7 до 10--11 років і характеризується інтенсивним розвитком довільності, рефлексії та здатності координувати свої дії з партнерами. Типова освітня програма для 3--4 класів інтегрованого курсу «Я досліджую світ» (за О.\,Я.~Савченко) прямо орієнтує педагогів на залучення учнів до дослідницьких практик, зокрема дослідження-розпізнавання, дослідження-спостереження і дослідження-пошуку, реалізація яких у груповому форматі відповідає природній соціальній потребі дітей цього вікового діапазону у спільній діяльності.

Поширене переконання про те, що першокласники занадто молоді для продуктивної групової роботи, спростовується емпіричними даними: цілеспрямоване дослідження поведінки семирічних учнів у малих групах із трьох осіб за методологією Lesson Study показало, що за відповідного педагогічного супроводу учні першого класу виявляють стійку здатність до кооперативної взаємодії, розподілу ролей та спільного конструювання результату \parencite{Jakavonytė-Staškuvienė2025}. Учні другого класу здатні виконувати inquiry-завдання з природознавства і технологій, розвиваючи при цьому наукову грамотність у суміщеному дослідницько-груповому форматі \parencite{Zuñiga-Quispe2025365}. На матеріалі четвертокласників встановлено нелінійний характер зв'язку між ступенем відкритості дослідницького завдання і рівнем навчальної автономії: широкий простір для inquiry одночасно провокує як найвищі, так і найнижчі рівні самостійності, що вимагає диференційованого підходу з боку вчителя \parencite{Aming'a2026}.

Соціально-емоційні умови групової взаємодії набувають особливого значення у молодшому шкільному віці. Дослідження управління «обличчям» (facework) у взаємодії учнів початкової школи виявило, що позитивна соціальна атмосфера у групі є необхідною умовою продуктивного пізнавального обміну між однолітками, тоді як загроза втрати «обличчя» суттєво пригнічує навчальний діалог \parencite{Evnitskaya2021169}. Аналіз групової динаміки четвертокласників під час колаборативного розв'язання математичних задач підтверджує, що відносини влади й статусу всередині групи є рухомими і потребують усвідомленого педагогічного управління \parencite{Langer-Osuna2020514}. Отже, організація групової взаємодії молодших школярів потребує врахування як вікових когнітивних можливостей, так і соціально-емоційного контексту спільної діяльності.

\newpage

% ==========================================================
% РОЗДІЛ 3 — ПІЗНАВАЛЬНА ДІЯЛЬНІСТЬ
% ==========================================================

\section{Пізнавальна діяльність у груповому форматі}

\subsection{Дослідницька поведінка учнів та inquiry-based learning}

Дослідницьке (inquiry-based) навчання як педагогічна стратегія ґрунтується на активному залученні учнів до процесів спостереження, формулювання запитань, висування гіпотез, збирання даних та осмислення результатів. Квазіекспериментальне дослідження ефективності inquiry-based learning (IBL) у змішаному синхронному навчальному середовищі підтвердило, що цей підхід суттєво підвищує залученість і навчальні результати учнів початкової школи порівняно з традиційним навчанням \parencite{Xu2026}. Якісний мета-аналіз міжнародних досліджень IBL засвідчив, що дослідницькі моделі навчання одночасно розвивають соціальні навички та покращують академічну успішність молодших школярів у різних культурних контекстах \parencite{Seprie2025}.

Дослідницька поведінка учнів початкової школи набуває різних форм залежно від навчального предмета і характеру завдання. Дослідження формування творчого мислення в учнів початкової школи на матеріалі природознавства підтвердило переваги дослідницько-орієнтованого підходу перед традиційним викладом \parencite{Leasa2025318}. Систематизація досвіду викладання природознавства і технологій у другому класі демонструє, що inquiry-підхід ефективно розвиває наукову грамотність вже у ранньому шкільному віці \parencite{Zuñiga-Quispe2025365}. Модель дослідницького навчання «Сім'я–Школа–Суспільство», розроблена в рамках китайської освітньої реформи, ілюструє потенціал розширення дослідницького середовища учнів за межі класної кімнати \parencite{Xia2025340}. Взаємодія учнів початкової школи з фізичними симуляціями у форматі дослідницького навчання сприяє осмисленню абстрактних понять через активне маніпулювання і перевірку гіпотез \parencite{Magkouta2024253}.

Окремий вектор досліджень пов'язаний із дослідницьким навчанням у контексті екологічної освіти та громадянської відповідальності. Якісне дослідження, присвячене формуванню «екологічного громадянства» у молодших школярів, показало, що inquiry-практики, спрямовані на вивчення місцевого біорізноманіття, розвивають як природничі компетентності, так і ціннісні орієнтації учнів щодо навколишнього середовища \parencite{Christodoulou2025969}. Тривале кейс-дослідження (13 місяців) засвідчило, що систематична дослідницька діяльність учнів 4--5 класів, опосередкована технологічним середовищем, формує стійкі пізнавальні стратегії та навички спільного конструювання знань \parencite{Viilo2018407}. Наведені дані дозволяють розглядати inquiry-based learning як методологічну основу пізнавальної діяльності учнів на заняттях «Я досліджую світ».


\subsection{Роль учителя-фасилітатора}

Перехід від моделі вчителя-транслятора знань до моделі вчителя-фасилітатора є однією з ключових характеристик сучасної педагогіки початкової школи, зокрема в контексті inquiry-based і кооперативного навчання. Порівняльне дослідження практик двох учителів початкової школи Південної Африки з різним педагогічним стажем показало, що досвідчені педагоги більш гнучко чергують пряму інструкцію і відкриті дослідницькі ситуації, тоді як початківці схильні до надмірного структурування завдань \parencite{Photo20251493}. Аналіз даних австрійського TIMSS-2019 дослідження виявив, що вчителі початкової школи, які частіше застосовують дослідницькі стратегії навчання, демонструють вищу чутливість до складу класу і гнучкіше адаптують інструкцію до потреб учнів \parencite{Neubauer-Hametner2025}.

Фасилітація спільного конструювання знань у цілому класі вимагає від учителя особливих комунікативних стратегій. Кейс-дослідження у чотирьох норвезьких початкових класах засвідчило, що перехід від закритих питань і оцінювальних реакцій до відкритих запитань і підтримки учнівських ідей суттєво підвищує рівень спільного конструювання знань у класній дискусії \parencite{Hogsnes20241}. Однорічна програма професійного розвитку підтвердила, що вчителі початкової школи здатні опанувати фасилітацію академічної аргументації та діалогічного обговорення текстів за умови систематичної методичної підтримки \parencite{Reznitskaya2019254}.

Технологічно опосередкована фасилітація відкриває нові можливості для організації inquiry-навчання. Експериментальне дослідження показало, що мобільні підказки-орієнтири під час дослідницьких завдань поза класною кімнатою підвищують якість дослідницького процесу і знижують когнітивне навантаження на учнів \parencite{Martin20245}. Взаємодія вчителя з фізичними симуляціями порівняно зі студентськими самостійними маніпуляціями виявляє переваги і обмеження кожної моделі \parencite{Magkouta2024253}. Водночас дослідження досвіду студентів педагогічних спеціальностей із inquiry-послідовностями для початкової школи виявляє, що майбутні вчителі часто недооцінюють складність відкритих форм дослідницького навчання, що актуалізує потребу в цілеспрямованій фаховій підготовці до ролі фасилітатора \parencite{Haslekås2022}.

\newpage

% ==========================================================
% РОЗДІЛ 4 — МЕТОДИЧНІ ПІДХОДИ
% ==========================================================

\section{Методичні підходи до організації групової роботи}

\subsection{Структурування завдань для групової дослідницької діяльності}

Якість групової дослідницької діяльності значною мірою визначається тим, як сформульоване і структуроване навчальне завдання. Концептуальна рамка проєктно-орієнтованого STEM-навчання (PjBL-STEM), розроблена для початкової школи, виокремлює кілька обов'язкових компонентів ефективного групового завдання: автентичність проблеми, міждисциплінарний характер, чіткий кінцевий продукт і вбудований механізм зворотного зв'язку \parencite{Retno2025579}. Реалізація цих принципів у практиці STEAM-проєктів початкової школи демонструє, що завдання, засновані на реальних суспільних проблемах (зокрема, кліматичних змінах і сталому розвитку міст), підвищують мотивацію учнів і забезпечують природну інтеграцію предметних змістів у спільній проєктній роботі \parencite{Vicente20201}.

Структура дослідницького групового завдання визначає характер когнітивної та соціальної взаємодії учасників. Дослідження реалізації трьох дизайн-орієнтованих STEM-проєктів у початкових класах показало, що завдання з елементами інженерного проєктування суттєво розвивають наукову творчість учнів і вимагають від групи узгодженого переходу між фазами ідеації, прототипування й оцінювання \parencite{Zhang2025568}. Розробка та апробація інтегрованих засобів навчання природознавства і суспільствознавства, орієнтованих на метакогнітивний розвиток учнів, засвідчила, що тематичне структурування міждисциплінарних завдань покращує критичне мислення молодших школярів \parencite{Sururuddin2026}. Аналіз кооперативного викладання STEAM у початковій школі за design-based підходом підтверджує, що спільна розробка та виконання проєктних завдань є ефективним механізмом інтеграції предметних знань і формування дослідницьких компетентностей \parencite{Li2022}.

Цифрові платформи відкривають нові можливості для структурування і підтримки групової дослідницької діяльності. Оцінювання педагогіко-технологічних характеристик платформ для проєктного навчання виявило, що найефективніші інструменти підтримують одночасно кілька фаз спільної роботи — від постановки завдання і розподілу ролей до документування дослідницького процесу і презентації результатів \parencite{Baziukė2025}. Дослідження впливу перцептомоторного досвіду на наукове дослідження молодших школярів доводить, що завдання, які залучають сенсорне сприйняття і фізичне маніпулювання, створюють міцнішу когнітивну основу для подальшої групової рефлексії \parencite{Farris2019648}. Наведені дані вказують на необхідність системного підходу до проєктування групових дослідницьких завдань, що враховує одночасно змістову, соціальну та інструментальну виміри спільної діяльності учнів.

% ============================================================
%  Розділ 4.2 — Оцінювання групової взаємодії
%  ~250 слів
% ============================================================

\subsection{Оцінювання групової взаємодії}

Оцінювання в контексті групової дослідницької діяльності є методично складним завданням, що потребує поєднання індивідуальних та групових показників. Дослідження впливу Програми оцінювання кооперативного навчання на соціальну взаємодію учнів 3--6 класів початкової школи Китаю (165 учасників) показало, що поєднання індивідуальної і групової складових оцінювання і фокус одночасно на академічних та соціальних результатах суттєво підвищує якість міжособистісної взаємодії в групах \parencite{Zeng2025317}. Ці дані узгоджуються з результатами більш раннього дослідження, у якому залучення учнів до розробки критеріїв взаємооцінювання з використанням мобільних пристроїв позитивно вплинуло на їхні досягнення та усвідомленість власних навчальних стратегій \parencite{Lai2015149}.

Взаємооцінювання є особливо ефективним інструментом у контексті дослідницького розв'язання задач. Якісний аналіз процесу розв'язання математичних задач шестикласниками засвідчив, що взаємооцінювання як складова нелінійного орієнтувального базису суттєво прискорює еволюцію дослідницьких стратегій групи і підвищує усвідомленість кожного учасника щодо якості власних міркувань \parencite{Torregrosa202189}. Рамка ECEE (дослідження~-- створення~-- вираження~-- оцінювання), розроблена для підготовки вчителів початкової школи, виокремлює оцінювання як окрему і рівнозначну фазу творчо-колаборативної діяльності, а не лише підсумковий захід \parencite{Mariati2026}.

Дослідження здатності майбутніх учителів початкової школи проєктувати дослідницькі завдання виявило систематичні прогалини у формулюванні вимірюваних критеріїв оцінювання, зокрема в ситуаціях із контрінтуїтивними результатами, що свідчить про потребу в цілеспрямованій підготовці педагогів до оцінювання дослідницького процесу, а не лише його продуктів \parencite{Lefkos2025349}. Оцінювання групових навчальних продуктів на цифрових навчальних платформах, у свою чергу, потребує врахування не лише академічних, а й колаборативних показників \parencite{Baziukė2025}.

% ============================================================
%  Розділ 5 — Інтеграція природничих, соціальних та громадянознавчих змістів
%  ~350 слів
% ============================================================

\newpage

% ==========================================================
% РОЗДІЛ 5 — ІНТЕГРАЦІЯ
% ==========================================================

\section{Інтеграція природничих, соціальних та громадянознавчих змістів
         у курсі «Я досліджую світ»}

Інтегрований курс «Я досліджую світ» у рамках Нової української школи поєднує природничу, соціальну, громадянознавчу та здоров'язбережувальну освітні галузі в єдиний навчальний простір, орієнтований на дослідницьку діяльність учнів. Типова освітня програма 3--4 класів (за О.\,Я.~Савченко) передбачає використання інтегрованих завдань дослідницького характеру, що охоплюють довкілля, суспільство та людину як взаємопов'язані об'єкти пізнання. Теоретичним обґрунтуванням такого підходу є ідея про те, що розподіл знань на ізольовані дисципліни суперечить природі дитячого мислення, тоді як інтеграція змістів дозволяє учням встановлювати міжпредметні зв'язки і бачити цілісність навколишнього світу \parencite{Lehane20261}.

Досвід реалізації інтегрованих підходів у початковій освіті різних країн підтверджує педагогічну обґрунтованість такого проєктування навчального змісту. Розробка і апробація тематичних засобів навчання, що інтегрують природничий і суспільствознавчий зміст у єдині дослідницькі завдання, суттєво покращує критичне мислення молодших школярів \parencite{Sururuddin2026}. Дослідження сприйняття вчителями початкової школи педагогіки глибинного навчання в умовах культурно-інтегрованої STEM-освіти засвідчило готовність педагогів до міждисциплінарного проєктування навчального процесу за наявності відповідної методичної підтримки \parencite{Putri20251331}.

STEM/STEAM-освіта є найбільш операціоналізованою формою інтеграції природничого та технологічного змістів на рівні початкової школи. Реалізація STEAM-проєктів з використанням освітньої роботики у п'ятому класі початкової школи демонструє, як кооперативне навчання і проєктна діяльність органічно поєднуються в єдиному міждисциплінарному форматі \parencite{Vicente20201}. Design-based дослідження впровадження STEAM-навчання через кооперативне викладання підтвердило, що міждисциплінарне конструювання знань формує в учнів початкової школи інформаційну грамотність і дослідницькі компетентності вищого порядку \parencite{Li2022}. Практичні дослідницькі активності у природознавстві суттєво підвищують інтерес молодших школярів до STEM-дисциплін, формуючи стійку мотивацію до наукового пізнання вже у початкових класах \parencite{Ribarić2025}. Застосування ICT-середовищ із високим рівнем візуалізації для вивчення STEM-понять у початковій школі знижує рівень абстракції навчального матеріалу і стимулює дослідницьку допитливість учнів \parencite{Fanchamps2023102}.

Проведений аналіз дозволяє стверджувати, що методична організація групової взаємодії учнів у пізнавальній діяльності на заняттях «Я досліджую світ» потребує узгодження трьох складових: інтегрованого змісту, дослідницьких методів і кооперативної форми організації навчання. Тільки системне поєднання цих компонентів забезпечує досягнення очікуваних результатів, задекларованих у типовій освітній програмі НУШ.

\newpage

% ==========================================================
% ВИСНОВКИ
% ==========================================================

\section*{Висновки огляду}
\addcontentsline{toc}{section}{Висновки огляду}

Проведений наративний огляд дозволяє сформулювати такі основні висновки. По-перше, кооперативне навчання в початковій школі є доказово ефективною стратегією, що одночасно розвиває академічні результати, соціально-емоційні навички та пізнавальну самостійність учнів, починаючи вже з першого класу. По-друге, inquiry-based learning забезпечує методологічну основу для організації дослідницької пізнавальної діяльності: відкритість завдань, автентичність проблем і залучення дітей до наукових практик є ключовими чинниками ефективності цього підходу в початкових класах. По-третє, роль учителя-фасилітатора є визначальною для якості групової взаємодії: комунікативні стратегії педагога, рівень його методичної готовності та здатність адаптувати ступінь підтримки до потреб учнів суттєво впливають на результативність спільної роботи. По-четверте, інтеграція природничих, соціальних та громадянознавчих змістів у дусі програми «Я досліджую світ» органічно поєднується з груповими дослідницькими практиками.

Виявлені прогалини в дослідженнях: обмаль робіт, що безпосередньо вивчають групову взаємодію в умовах інтегрованих курсів типу «Я досліджую світ»; брак досліджень у контексті НУШ та пострадянських освітніх систем; недостатнє охоплення вікової групи 3--4 класів як самостійного об'єкта вивчення. Ці прогалини визначають актуальність і новизну подальшого дослідження в курсовій роботі.

\newpage

% ==========================================================
% СПИСОК ВИКОРИСТАНИХ ДЖЕРЕЛ
% ==========================================================

\printbibliography[title={Список використаних джерел}]

\end{document}
